% Chapter 2: a surface-level theme vs. a theme that goes one level deeper.
% Shared preamble for the book's diagrams.
%
% Each figures/<name>.tex is a standalone TikZ picture. `make figures` builds,
% for each language <lang> (ko, en),
%   img/<lang>/<name>.pdf  vector, used by the print book (book.lua swaps it in)
%   img/<lang>/<name>.png  300dpi, white background, used by the Markdown
%
% Text is written once per language as \tr{한국어}{English}. The language is
% chosen by defining \figlang before this file is read, e.g.
%   xelatex -jobname=<name> "\def\figlang{en}\input{figures/<name>}"
% and defaults to ko.
%
% Figures are drawn to fit the print text block (95mm wide) at 1:1 scale, in
% grayscale only so they survive black-and-white printing.
\documentclass{article}
\providecommand\figlang{ko}
\def\figlangKo{ko}
\ifx\figlang\figlangKo
  \newcommand\tr[2]{#1}
\else
  \newcommand\tr[2]{#2}
\fi
\usepackage{fontspec}
\ifx\figlang\figlangKo
  \usepackage[hangul]{xetexko}
  \setmainfont{Noto Sans CJK KR}
  \setmainhangulfont{Noto Sans CJK KR}
\else
  \setmainfont{Noto Sans}
  % No hyphenation: diagram labels read better with whole words.
  \AtBeginDocument{\hyphenpenalty=10000 \exhyphenpenalty=10000 }
\fi
\usepackage{xcolor}
\usepackage{tikz}
\usetikzlibrary{arrows.meta,positioning,calc,fit,decorations.pathmorphing,decorations.pathreplacing,shapes.geometric,backgrounds}

% Grays (print-safe)
\definecolor{ink}{gray}{0.12}   % lines and text
\definecolor{mid}{gray}{0.45}   % secondary text, guides
\definecolor{soft}{gray}{0.72}  % light strokes
\definecolor{lite}{gray}{0.90}  % filled shapes (sticky notes, boxes)
\definecolor{pale}{gray}{0.96}  % large background areas

% Font sizes in pt, relative to the 9.5pt body text.
\newcommand\figtext{\fontsize{7.5}{10}\selectfont}
\newcommand\figsmall{\fontsize{6.5}{8.5}\selectfont}
\newcommand\figtiny{\fontsize{5.5}{7}\selectfont}

\tikzset{
  every picture/.style={line width=0.6pt, draw=ink, text=ink, font=\figtext, >={Stealth[length=2.2mm,width=1.6mm]}},
  arr/.style={->, draw=ink, line width=0.6pt},
}
\setlength\parindent{0pt}
\pagestyle{empty}
% Break Korean lines only at spaces (between words), not between syllables,
% which reads better in narrow diagram boxes.
\ifx\figlang\figlangKo
  \AtBeginDocument{\XeTeXlinebreakpenalty=10000 }
\fi

% Crop each page to its picture (+1mm border); TeX Live here has no
% standalone/preview packages. Wrap the tikzpicture in \begin{fig}...\end{fig}.
\newsavebox\figbox
\newenvironment{fig}{\begin{lrbox}{\figbox}}{\end{lrbox}\shipfig}
\newcommand\shipfig{%
  \hoffset=-1in \voffset=-1in
  \pdfpagewidth=\dimexpr\wd\figbox+2mm\relax
  \pdfpageheight=\dimexpr\ht\figbox+\dp\figbox+2mm\relax
  \shipout\vbox{\kern1mm\hbox{\kern1mm\usebox\figbox}\kern\dp\figbox\kern1mm}}

\begin{document}\begin{fig}
\begin{tikzpicture}[
    card/.style={text width=84mm, inner sep=2.5mm, align=left, font=\figsmall},
    head/.style={font=\figsmall\bfseries, anchor=south west, inner sep=0.5mm},
  ]
  \node[card, draw=soft, dashed, line width=0.6pt, anchor=north west] (bad) at (0,0)
    {\tr{음료를 기다리는 고객들이 카페 입구에서 대기하기 때문에 새롭게 입장하는 고객들이 불편을 느끼고 입장 후 주문까지 시간이 오래 걸린다.}%
         {Customers waiting for their drinks stand at the café entrance, so people coming in are inconvenienced and it takes a long time from entering to ordering.}\par
     {\color{mid}\tr{→}{$\rightarrow$} \tr{현상을 묘사하는 데 그친다.}{It only describes what happens.}}};
  \node[head, text=mid] at ($(bad.north west)+(0,0.8mm)$) {\tr{표면적인 주제}{Surface-level theme}};

  \node[card, draw=ink, line width=1.1pt, anchor=north west] (good) at ($(bad.south west)+(0,-11mm)$)
    {\tr{주문하는 곳과 음료를 받는 곳이 나란히 붙어 있기 때문에 음료를 대기하는 고객들이 입장 고객의 동선을 막게 되고 모든 고객들이 불편을 느낀다.}%
         {The order counter and the pickup counter sit right next to each other, so customers waiting for drinks block the way of people coming in, and every customer is inconvenienced.}\par
     {\color{mid}\tr{→}{$\rightarrow$} \tr{키오스크는 지금 주문을 받는 곳과 다른 위치에 있어야 한다는 함의가 보인다.}{This suggests the kiosk should stand somewhere other than where orders are taken now.}}};
  \node[head] at ($(good.north west)+(0,0.8mm)$) {\tr{한 단계 더 들어간 주제}{A theme that goes one level deeper}};

  \draw[arr, draw=mid] ($(bad.south)+(0,-0.5mm)$) -- node[right=1mm, font=\figsmall, text=mid] {\tr{왜 입구에서 기다릴 수밖에 없을까?}{Why must they wait at the entrance?}} ($(good.north)+(0,4.5mm)$);
\end{tikzpicture}
\end{fig}\end{document}
